\section{Threats to validity\label{sec:threats}}

\paragraph{Construct validity.}
\emph{Junior}, \emph{Mid}, and \emph{Senior} are unsupervised clusters over observable GitHub activity, not verified job titles or measured skill; a developer's placement reflects relative activity within the analyzed population, not a career stage. Three of the twelve clustering indicators -- Agentic-PR activity span, agent diversity, and Agentic-PR repository breadth -- directly describe agent-associated activity rather than independent, pre-agent experience, so the profiles are partly constructed from the same kind of behavior RQ1 and RQ2 later measure; the remaining nine (GraphQL-derived contribution history, plus followers, account tenure, and repository stars/forks) describe general GitHub activity not tied to agent use. The RQ3 outcome counts non-bot comments from any other account, not comments verified to come from a maintainer; it is best read as a proxy for visible discussion volume, not review effort or defect count. Bot and agent detection relies on AIDev's own labeling and simple login-pattern filtering, which may miscount atypically named bot or human accounts. The GraphQL profile enrichment was collected on 6 May 2026, about nine months after AIDev's own 1 August 2025 cutoff; indicators such as followers, account tenure, and yearly contribution counts reflect each account's state at the later collection date, not at the time of its Agentic-PRs, so part of the activity profile may postdate the contributions it is meant to characterize.

\paragraph{Internal validity.}
The design is observational: none of the four research questions supports a causal claim about the effect of experience or of agent use. PR complexity, task type, and repository governance norms are not controlled for and plausibly confound both acceptance rate and comment volume; RQ2's finding that acceptance decreases with activity level, for instance, is also consistent with more active developers attempting more difficult tasks rather than with agents serving them less well. The $k=3$ operationalization was retained for interpretability even though neither silhouette width nor the gap statistic identifies it as the strict numerical optimum (Section~\ref{sec:method}); a different $k$ or a different clustering algorithm could yield different group compositions and, in turn, different outcome comparisons. The 1,755 developers with GraphQL-imputed indicators were estimated from AIDev-derived features with weak explanatory power (observed $R^2$ between $-0.27$ and $0.10$, Section~\ref{sec:method}); their placement into activity profiles is comparatively uncertain.

A complete-case sensitivity check -- clustering only the 70,434 developers with real, non-imputed GraphQL data, instead of the full 72,189 -- initially found a quantile-transformation-based version of the pipeline in Section~\ref{sec:method} to be unstable under this population change: cluster assignment agreed with the full-population clustering, for the same overlapping developers, at only Adjusted Rand Index $=0.46$, and RQ1's and RQ2's Junior-vs-Mid ordering reversed. We traced this to the quantile transformation step: because several indicators are heavily zero-inflated, quantile transformation maps values by rank within the current sample, so removing even a small fraction of the population shifts those ranks for everyone else. Re-running the same complete-case comparison with $\log(1+x)$ transformation followed by standardization in place of quantile transformation gave Adjusted Rand Index $=0.98$, essentially stable, which is why $\log(1+x)$ plus standardization -- not quantile transformation -- is the preprocessing step used throughout this paper (Section~\ref{sec:method}).

We also reran RQ1--RQ3 restricted to the 70,434 complete-case developers under the now-adopted $\log(1+x)$ pipeline (RQ4 does not depend on clustering and is unaffected). RQ1 and RQ2 were essentially unchanged (e.g.\ RQ1: Junior/Mid/Senior means $3.59/42.27/7.18$ vs.\ $3.57/42.12/7.94$ in the full-population analysis; RQ2: $90.3\%/91.8\%/83.4\%$ vs.\ $90.2\%/91.8\%/83.2\%$), with the same pairwise significance pattern in both cases. \textbf{RQ3 did not replicate}: under the complete-case restriction its already-small AIDev-pop sample shrinks further and shifts composition (Senior drops from 13,434 to 4,366 observations), the mean ordering reverses (Junior $0.45$, Mid $0.56$, Senior $0.73$, i.e.\ increasing rather than decreasing with activity level), and the Kruskal--Wallis test is no longer significant ($p=0.77$). RQ3 is therefore the least robust of the four research questions: its direction and significance both depend on whether the 1,755 imputed developers are included, on top of the small-sample concerns already noted below.

We additionally ran four further sensitivity checks against the $\log(1+x)$-based pipeline. (i) \emph{Alternative $k$}: at $k=4$ and $k=5$, the three main clusters remain close to their $k=3$ sizes, with the additional clusters carved out of Senior as small, higher-activity splinters, rather than a wholesale reorganization. (ii) \emph{Repeated initializations}: rerunning the full pipeline (Isolation Forest and $k$-means) with five different random seeds gave Adjusted Rand Index between $0.98$ and $0.99$ against the reported clustering -- highly stable. (iii) \emph{Retained outliers}: skipping Isolation Forest entirely gives Adjusted Rand Index $=0.64$ against the reported clustering, driven almost entirely by the Senior profile, which shrinks from 13,908 to 1,586 developers as the most extreme profiles get isolated into their own cluster instead of being excluded beforehand; Junior and Mid are comparatively stable. (iv) \emph{Excluding the AIDev-derived indicators} (Agentic-PR activity span, agent diversity, and repository breadth), leaving only the nine remaining indicators: this is the least stable of the four checks, Adjusted Rand Index $=0.32$, and cluster sizes shift substantially (Mid grows from 12,705 to 28,329). We re-ran RQ1 and RQ2 under this 9-indicator clustering: RQ2's qualitative finding holds -- Senior still has the lowest acceptance rate (80.7\% vs.\ 90.4\% Junior and 90.2\% Mid) -- but RQ1's Mid-vs-Senior contrast is much weaker (Agentic-PR means 17.6 vs.\ 20.8, close to reversing) than under the full 12-indicator clustering, consistent with the construct-validity concern raised above: part of RQ1's large effect size is attributable to the clustering itself being partly built from agent-usage indicators.

\paragraph{External validity.}
The dataset covers Agentic-PRs from five specific coding agents on public GitHub repositories up to 1 August 2025, an early period of agent adoption; usage patterns, acceptance norms, and maintainer behavior around agent-generated code may shift as the tooling and community norms mature, so the reported associations should not be extrapolated to later periods or to private/enterprise repositories without re-verification. RQ3 and RQ4 are further restricted to AIDev-pop and to a $>500$-star baseline respectively, because richer interaction data (comments, reviews, commits, issues) and the human-authored comparison sample were collected by AIDev's authors only for repositories above those popularity thresholds; both results describe high-visibility, high-star repositories specifically and may not generalize to the long tail of smaller projects that make up most of the full-scope population.

\paragraph{Conclusion validity.}
The outcome distributions are strongly non-normal and zero-inflated, which motivated the nonparametric tests and the Bonferroni correction used throughout (Section~\ref{sec:method}). With sample sizes in the tens of thousands, very small differences can reach statistical significance regardless of practical importance; this applies most clearly to RQ2 ($\varepsilon^2=0.043$) and especially RQ3 ($\varepsilon^2=0.009$), whose effect sizes are small even though their $p$-values are not. RQ1 is an exception: its effect size ($\varepsilon^2=0.32$) is large by conventional standards, so its significance reflects a substantial association, not only a large sample. RQ3 and RQ4 additionally rest on markedly unequal group sizes (RQ3's Junior profile has only 187 observations against 13,434 for Senior; RQ4 compares 56,387 Agentic-PRs to 2,286 human-authored PRs), which further limits how much weight their point estimates, particularly for the smaller group, can bear.
